\section{Patch positivity and the critical profile}
\label{sec:patch-positivity}

\subsection{Definition and coefficient characterization}
\label{subsec:positivity-criterion}

\begin{definition}[Patch positivity]
\label{def:patch-positivity}
A spin system has \emph{patch positivity} if, for every~$A\Subset\Lambda$,
\begin{equation*}
C(P)\geq0
\qquad
\text{for every }P\in\mathcal P,
\quad \pP_A\text{-almost surely}.
\end{equation*}
\end{definition}

The end contributions depend additionally on the terminal configuration and
require a one-density threshold. Write
\begin{equation*}
r_i=c_i^0(\vn)+c_i^1(\vn).
\end{equation*}

\begin{theorem}[Coefficient criterion]
\label{thm:patch-positivity-criterion}
A spin system has patch positivity if and only if, for every~$i\in\Lambda$ and
nonempty~$S\subseteq N(i)$,
\begin{equation*}
c_i^0(S)+c_i^1(S)\leq0
\end{equation*}
and
\begin{equation*}
\pw{
 c_i^1(\vn)c_i^0(S)-c_i^0(\vn)c_i^1(S)\geq0
}{r_i>0}
{
 c_i^1(S)\leq0
}{r_i=0}.
\end{equation*}
\end{theorem}

\begin{proof}
Fix a site~$i$ and use the full-patch formulas from
Appendix~\ref{app:bulk-formulas}. For patches of type~$\mathsf{II}$ or
$\mathsf{IE}$,
\begin{equation*}
C(P)=\frac{\psi_i(\Delta,1)}{\varphi_i(\Delta)}\geq0.
\end{equation*}
 Patches of
 type~$\mathsf{IO}$ have contribution~$e^{V_i\Delta}>0$. These types impose no
condition on the nonconstant rate coefficients.

Let~$P$ start with an outgoing interaction of nonempty target~$S$. For a patch
of type~$\mathsf{OO}$,
\begin{equation*}
C(P)=\sigma_i^\beta(S)e^{V_i\Delta}.
\end{equation*}
Since
\begin{equation*}
\beta_i(S)\sigma_i^\beta(S)
=-c_i^0(S)-c_i^1(S),
\end{equation*}
these contributions are nonnegative exactly when
$c_i^0(S)+c_i^1(S)\leq0$.

It remains to consider types~$\mathsf{OI}$ and~$\mathsf{OE}$. Their denominator
is positive on every realized patch, and their numerator is
\begin{equation*}
c_i^0(S)
-\bb{c_i^0(S)+c_i^1(S)}\psi_i(\Delta,1).
\end{equation*}
Assume the condition obtained from the~$\mathsf{OO}$ patches. If~$r_i>0$, then
\begin{equation*}
\psi_i(\Delta,1)
=
\frac{c_i^0(\vn)}{r_i}
+
\frac{c_i^1(\vn)}{r_i}e^{-r_i\Delta}.
\end{equation*}
The empty coefficients,~$c_i^x(\vn) = c_i\bb{\mb{1}^{i,x}}$ are nonnegative, so the numerator is minimized in the long-patch limit, where it equals
\begin{equation*}
\frac{
 c_i^1(\vn)c_i^0(S)-c_i^0(\vn)c_i^1(S)
}{r_i}.
\end{equation*}
It is nonnegative for every patch length exactly when the quantity in the
statement is nonnegative. If~$r_i=0$, then both empty-neighbour flip rates
vanish and~$\psi_i(\Delta,1)=1$. The numerator reduces to~$-c_i^1(S)$, giving
the second case. This checks every patch type and proves the criterion.
\end{proof}

\subsection{The calm--facilitating consequence}
\label{subsec:positivity-interpretation}

\begin{corollary}
\label{cor:calm-facilitating-coefficients}
{\color{red} Isnt the more relevant statement that patch positivity makes $c_i$ nonincreasing when calm states are added to neighbourhood of i? Exception is case $r_i = 0$ but that one is trivial since it means no interactions (since $c_i^x = c_i^x(\vn)$ + negative terms$\geq 0$ so if $c_i^x(\vn)=0$ then the whole rate is zero). Thats a much better interpretation than this weird $c_i^1 + c_i^0$ rate that has no obvious interpretation. More facilitators -> more activity is the right interpretation. }
If a spin system has patch positivity, then for every~$i\in\Lambda$ and every
nonempty~$S\subseteq N(i)$,
\begin{equation*}
c_i^1(S)\leq0,
\qquad
c_i^0(S)+c_i^1(S)\leq0.
\end{equation*}
Consequently, the activation rate~$c_i^1$ and the total flip rate
$c_i^0+c_i^1$ are coordinatewise nonincreasing as calm states are added to the
neighbourhood of~$i$.
\end{corollary}

\begin{proof}
The second inequality is part of Theorem~\ref{thm:patch-positivity-criterion}.
When~$r_i=0$, the first one is also part of the criterion. Suppose that~$r_i>0$
and~$c_i^1(S)>0$. The second inequality then gives~$c_i^0(S)<0$, and hence
\begin{equation*}
c_i^1(\vn)c_i^0(S)-c_i^0(\vn)c_i^1(S)<0,
\end{equation*}
because~$c_i^0(\vn),c_i^1(\vn)\geq0$ and at least one is positive. This
contradicts the second condition. The coordinatewise statement follows
from the multilinear expansions: every nonconstant coefficient of~$c_i^1$ and
of~$c_i^0+c_i^1$ is nonpositive.
\end{proof}

This consequence formalizes the calm--facilitating tendency used in the
introduction. It is weaker than patch positivity. Coordinatewise decrease does
not in general force all multilinear coefficients to be nonpositive, and even
coefficientwise decrease does not impose the determinant {\color{red} weird phrasing, what determinant?} relation between the
two flip directions.

\subsection{End contributions and the critical profile}
\label{subsec:critical-density}

Patch positivity does not by itself determine the sign of an end contribution,
because~$C(p,P)$ also contains the terminal one-density~$p$. For a patch-positive
spin system, define the \emph{patch critical density} at~$i$ by
\begin{equation*}
p_i^\star
=
\inf\cb{
 p\in[0,1]:
 C(p,P)\geq0
 \text{ for every end patch }P\text{ based at }i
}.
\end{equation*}
The patch critical profile is~$\mb p^\star=(p_i^\star)_{i\in\Lambda}$. For a
translation-invariant system it is constant, and its common value is denoted
by~$p^\star$.

\begin{proposition}[Critical-profile formula]
\label{prop:critical-profile-formula}
For every~$i\in\Lambda$,
\begin{equation*}
p_i^\star
=
\max\cb{
0,
\sup_{\substack{
 \vn\neq S\subseteq N(i)\\
 c_i^0(S)+c_i^1(S)<0
}}
\frac{c_i^0(S)}{c_i^0(S)+c_i^1(S)}
}.
\end{equation*}
{\color{red} Unnecessary comment, the supremum of empty set is canonically $-\infty$.}
The inner supremum is~$0$ when its index set is empty.
\end{proposition}

\begin{proof}
An incoming end patch has contribution
$\psi_i(\Delta,p)/\varphi_i(\Delta)$ and is nonnegative for every
$p\in[0,1]$. Consider an outgoing end patch based at~$i$, with initial target
$S\neq\vn$. Its denominator is positive, and its numerator is
\begin{equation*}
N_{i,S}(\Delta,p)
=
c_i^0(S)
-\bb{c_i^0(S)+c_i^1(S)}\psi_i(\Delta,p).
\end{equation*}
Patch positivity gives~$c_i^0(S)+c_i^1(S)\leq0$, so this numerator is
nondecreasing in~$p$.
{\color{red} These constant $r_i = 0$ cases are dumb. Given patch positivity if $r_i=0$ then~$c_i = 0$ so there are no patches that have an outgoing interaction at site~$i$.  }
Suppose first that~$r_i>0$. Since
\begin{equation*}
\psi_i(\Delta,p)
=
\frac{c_i^0(\vn)}{r_i}
+
\bb{p-\frac{c_i^0(\vn)}{r_i}}e^{-r_i\Delta},
\end{equation*}
$N_{i,S}(\Delta,p)$ interpolates between its short-patch limit
\begin{equation*}
c_i^0(S)-\bb{c_i^0(S)+c_i^1(S)}p
\end{equation*}
and its long-patch limit
\begin{equation*}
\frac{
 c_i^1(\vn)c_i^0(S)-c_i^0(\vn)c_i^1(S)
}{r_i}.
\end{equation*}
The latter is nonnegative by Theorem~\ref{thm:patch-positivity-criterion}.
Therefore the contribution is nonnegative for every patch length exactly when
the short-patch limit is nonnegative. If~$r_i=0$, then
$\psi_i(\Delta,p)=p$, and the same condition is obtained directly.

When~$c_i^0(S)+c_i^1(S)=0$, the coefficient criterion gives
$c_i^0(S)\geq0$, so the target~$S$ imposes no restriction. When the sum is
negative, the short-patch condition is equivalent to
\begin{equation*}
p\geq
\frac{c_i^0(S)}{c_i^0(S)+c_i^1(S)}.
\end{equation*}
Taking the supremum over nonempty targets and imposing~$p\geq0$ proves the
formula. Patch positivity also ensures that every ratio entering the supremum
is at most~$1$.
\end{proof}

\begin{corollary}[Empty-neighbour bound]
\label{cor:critical-profile-empty-bound}
If~$r_i>0$, then
\begin{equation*}
p_i^\star
\leq
\frac{c_i^0(\vn)}{c_i^0(\vn)+c_i^1(\vn)}.
\end{equation*}
\end{corollary}

\begin{proof}
For every nonempty~$S\subseteq N(i)$ with
$c_i^0(S)+c_i^1(S)<0$, the determinant condition is equivalent to
\begin{equation*}
\frac{c_i^0(S)}{c_i^0(S)+c_i^1(S)}
\leq
\frac{c_i^0(\vn)}{c_i^0(\vn)+c_i^1(\vn)}.
\end{equation*}
The claim follows from Proposition~\ref{prop:critical-profile-formula}.
\end{proof}

For every end patch~$P$ based at~$i$, the function~$p\mapsto C(p,P)$ is affine
and nondecreasing, and it is nonnegative for~$p\geq p_i^\star$.
